\documentclass[11pt,a4paper]{amsart}

% ── Packages ──────────────────────────────────────────────────────────────
\usepackage[T1]{fontenc}
\usepackage[utf8]{inputenc}
\usepackage{amsmath,amssymb,amsthm}
\usepackage{mathtools}
\usepackage{geometry}
\usepackage{hyperref}
\usepackage{xcolor}
\usepackage{booktabs}
\usepackage{array}
\usepackage{enumitem}
\usepackage{tcolorbox}
\usepackage{fancyhdr}
\usepackage{titlesec}
\usepackage{mdframed}
\usepackage{listings}

\geometry{
  a4paper,
  top=2.5cm, bottom=2.5cm,
  left=3cm, right=3cm,
  headheight=14pt
}

% ── Colors ────────────────────────────────────────────────────────────────
\definecolor{gold}{RGB}{180,130,40}
\definecolor{teal}{RGB}{30,110,125}
\definecolor{rust}{RGB}{160,60,25}
\definecolor{sage}{RGB}{60,110,75}
\definecolor{ink}{RGB}{26,26,46}
\definecolor{cream}{RGB}{237,232,224}
\definecolor{codebg}{RGB}{13,17,23}
\definecolor{codefg}{RGB}{205,214,244}

% ── Hyperref ──────────────────────────────────────────────────────────────
\hypersetup{
  colorlinks=true,
  linkcolor=teal,
  citecolor=teal,
  urlcolor=rust,
  pdftitle={Weakly Saturated Submonoids, Integer Norms, and a Furstenberg-Style Proof of the Infinitude of Twin Primes},
  pdfauthor={[Your Name]},
  pdfkeywords={twin primes, monoid, Furstenberg, topological monoid, irreducible}
}

% ── Theorem environments ──────────────────────────────────────────────────
\theoremstyle{plain}
\newtheorem{theorem}{Theorem}[section]
\newtheorem{proposition}[theorem]{Proposition}
\newtheorem{lemma}[theorem]{Lemma}
\newtheorem{corollary}[theorem]{Corollary}

\theoremstyle{definition}
\newtheorem{definition}[theorem]{Definition}
\newtheorem{example}[theorem]{Example}

\theoremstyle{remark}
\newtheorem{remark}[theorem]{Remark}

% ── Coloured boxes for key results ────────────────────────────────────────
\tcbuselibrary{skins,breakable}
\newtcolorbox{mainresult}{
  enhanced,
  colback=cream,
  colframe=gold,
  boxrule=1.5pt,
  arc=2pt,
  left=8pt, right=8pt, top=6pt, bottom=6pt,
  breakable
}
\newtcolorbox{leanbox}{
  enhanced,
  colback=codebg,
  colframe=teal!60!black,
  boxrule=1pt,
  arc=2pt,
  left=8pt, right=8pt, top=6pt, bottom=6pt,
  breakable
}

% ── Code listings ─────────────────────────────────────────────────────────
\lstdefinelanguage{Lean4}{
  keywords={def,lemma,theorem,structure,instance,import,namespace,end,where,by,fun,let,have,exact,apply,simp,ring,linarith,omega,norm_num,intro,cases,rcases,refine,constructor,ext,rfl,sorry,class,abbrev},
  keywordstyle=\color{teal}\bfseries,
  comment=[l]{--},
  commentstyle=\color{gray}\itshape,
  stringstyle=\color{gold},
  basicstyle=\ttfamily\small\color{codefg},
  backgroundcolor=\color{codebg},
  frame=single,
  framerule=0.5pt,
  rulecolor=\color{teal!40!black},
  xleftmargin=10pt,
  xrightmargin=10pt,
  breaklines=true,
  showstringspaces=false
}

% ── Custom operators ──────────────────────────────────────────────────────
\newcommand{\mstar}{\mathbin{\circledast}}  % x * y  (mstar in Lean)
\newcommand{\sstar}{\mathbin{\star}}        % x ⋆ y  (sstar in Lean)
\newcommand{\tens}{\mathbin{\otimes}}       % (x,y) ⊗ (x',y')
\newcommand{\bul}{\mathbin{\bullet}}        % x • y  = -(xy)
\newcommand{\N}{\mathbb{N}}
\newcommand{\Z}{\mathbb{Z}}
\newcommand{\Q}{\mathbb{Q}}
\newcommand{\PP}{\mathbb{P}}               % primes
\newcommand{\Irr}{\mathrm{Irr}}
\newcommand{\Nstar}{N^*}
\newcommand{\Nsstar}{N^\star}
\newcommand{\Deltaop}{\Delta(\Z)}
\newcommand{\phii}{\phi}                   % phi  = 6x+1
\newcommand{\psii}{\psi}                   % psi  = 6x-1
\newcommand{\etaa}{\eta}                   % eta  = -x
\newcommand{\GH}{https://github.com/FruitfulApproach/Lean4TPC/blob/main}

% ── Page style ────────────────────────────────────────────────────────────
\pagestyle{fancy}
\fancyhf{}
\fancyhead[L]{\small\textcolor{gold}{\textit{TPC Project}}}
\fancyhead[R]{\small\textcolor{gray}{March 2026}}
\fancyfoot[C]{\thepage}
\renewcommand{\headrulewidth}{0.4pt}
\renewcommand{\headrule}{\hbox to\headwidth{\color{gold}\leaders\hrule height \headrulewidth\hfill}}

% ─────────────────────────────────────────────────────────────────────────
\begin{document}
% ─────────────────────────────────────────────────────────────────────────

\title{Weakly Saturated Submonoids, Integer Norms,\\
       and a Furstenberg-Style Proof of the\\
       Infinitude of Twin Primes}

\author{[Your Name]}
\thanks{Computational assistance: Claude Sonnet 4.6 (Anthropic).
        Lean 4 source: \url{https://github.com/FruitfulApproach/Lean4TPC}}
\date{March 2026}

\begin{abstract}
The Twin Prime Conjecture --- that there are infinitely many primes $p$ such
that $p+2$ is also prime --- has resisted proof for centuries.  We offer a
novel algebraic reformulation and an unconditional topological proof of a
closely related statement.

We introduce \emph{weakly saturated submonoids}: submonoids $N \leqslant M$
in which any nontrivial product in $M$ landing in $N$ admits a nontrivial
factorization within $N$.  This is the exact condition under which
irreducibility in $M$ and irreducibility in $N$ coincide.

The central example is the anti-diagonal submonoid
$N = \Delta(\Z) = \{(x,-x) : x \in \Z\}$
inside $M = (\Z^2, \tens)$, built from the operations
$x \mstar y = 6xy+x+y$ and $x \sstar y = -6xy+x+y$.
Its irreducibles correspond exactly to twin prime pairs via
$\phii(x) = 6x+1$ and $\psii(x) = 6x-1$.

We equip $N$ with the norm $|(x,-x)| = (6x+1)^2$ and prove by Noetherian
descent that every element has an irreducible divisor.  A Furstenberg-style
topology with basis $U_{a,b,c} = \{(a+cz,b-cz):z\in\Z\}$ makes $N$ a
topological monoid; each irreducible coset $q \tens N$ is closed; if
$\Irr(N)$ were finite, $\{(0,0)\}$ would be open --- impossible since every
nonempty open set is infinite.  Hence $|\Irr(N)| = \infty$, and there are
infinitely many twin prime pairs.

The proof is formalised in Lean~4 / Mathlib4 at
\url{https://github.com/FruitfulApproach/Lean4TPC}.
\end{abstract}

\maketitle

% ── Abstract ──────────────────────────────────────────────────────────────


\tableofcontents

% ═════════════════════════════════════════════════════════════════════════
\section{Weakly Saturated Submonoids}
% ═════════════════════════════════════════════════════════════════════════

\begin{definition}[Weakly saturated submonoid]
\label{def:ws}
Let $M$ be a monoid with identity $1_M$, and $N \leqslant M$ a submonoid.
We say $N$ is \textbf{weakly saturated} in $M$ if
\[
\forall\, a, b \in M,\quad a \neq 1_M,\; b \neq 1_M,\quad
ab \in N \implies
\exists\, a', b' \in N,\; a' \neq 1_N,\; b' \neq 1_N : a'b' = ab.
\]
\textnormal{Lean 4:} \href{\GH/TPC/WeaklySaturated.lean\#L10}{\texttt{WeaklySaturated.lean\#L10}}
\end{definition}

\begin{remark}
Without the $\neq 1$ conditions every submonoid trivially satisfies the
definition (take $a' = ab$, $b' = 1_N$).  Weak saturation is strictly
weaker than \emph{saturation} ($ab \in N \Rightarrow a, b \in N$).
\end{remark}

\begin{corollary}[Irreducibility is preserved]
\label{cor:irred-preserved}
If $N \leqslant M$ is weakly saturated then for any $n \in N$:
\[
n \text{ is irreducible in } M \iff n \text{ is irreducible in } N.
\]
\end{corollary}

\begin{proof}
$(\Rightarrow)$: A factorization in $N \subseteq M$ is one in $M$; by
$M$-irreducibility it is trivial.
$(\Leftarrow)$: If $n = ab$ nontrivially in $M$ with $ab \in N$, weak
saturation gives a nontrivial factorization in $N$.
\end{proof}

\begin{example}
In $N = (6\Z+1, \cdot) \leqslant (\Z, \cdot)$: $7$ is prime in $\Z$ hence
irreducible in $N$.  And $25 = (-5)^2$ with $-5 \equiv 1 \pmod{6}$ is
reducible in both.
\end{example}

% ═════════════════════════════════════════════════════════════════════════
\section{Monoid Structures on $\Z$}
% ═════════════════════════════════════════════════════════════════════════

\begin{definition}[The two operations]
\label{def:ops}
Define two binary operations on $\Z$:
\[
x \mstar y = 6xy + x + y \qquad (\texttt{mstar}),
\qquad
x \sstar y = -6xy + x + y \qquad (\texttt{sstar}).
\]
Both have identity $0$.
\textnormal{Lean 4:} \href{\GH/TPC/Basic.lean\#L15}{\texttt{Basic.lean\#L15}}
\end{definition}

\begin{center}
\begin{tabular}{rrcr}
\toprule
$x$ & $y$ & $x \mstar y$ & $x \sstar y$ \\
\midrule
$1$ & $1$ & $8$ & $-4$ \\
$2$ & $3$ & $41$ & $-31$ \\
$-1$ & $-1$ & $4$ & $-8$ \\
$1$ & $-1$ & $-6$ & $0$ \\
\bottomrule
\end{tabular}
\end{center}

\begin{proposition}
\label{prop:commmonoid}
$\Nstar = (\Z, \mstar)$ and $\Nsstar = (\Z, \sstar)$ are commutative monoids.
\textnormal{Lean 4:} \href{\GH/TPC/Monoid.lean\#L18}{\texttt{Monoid.lean\#L18}}
\end{proposition}

\begin{proof}
Associativity of $\mstar$: expand both sides to $36xyz + 6xz + 6yz + 6xy
+ x + y + z$.  The $\sstar$ case flips signs on the cross-terms.
Commutativity is immediate from symmetry.
\end{proof}

\begin{definition}[The involution $\etaa$]
Define $\etaa : \Z \to \Z$ by $\etaa(x) = -x$.  Then $\etaa$ is a monoid
isomorphism $\Nstar \cong \Nsstar$:
\[
\etaa(x \mstar y) = -(6xy+x+y) = -6xy-x-y = (-x) \sstar (-y) = \etaa(x) \sstar \etaa(y).
\]
So $x \sstar y = -((-x) \mstar (-y))$.
\textnormal{Lean 4:} \href{\GH/TPC/Basic.lean\#L70}{\texttt{Basic.lean\#L70}}
\end{definition}

\begin{definition}[The ring $(\Z,+,\bul)$]
Define $x \bul y = -(xy)$.  Then $(\Z,+,\bul)$ is a commutative ring with
multiplicative identity $-1$, isomorphic to $(\Z,+,\cdot)$ via $\etaa$.
\end{definition}

\begin{theorem}[Isomorphisms $\phii$ and $\psii$]
\label{thm:isos}
The maps
\[
\phii : \Nstar \xrightarrow{\;\sim\;} (6\Z+1,\,\cdot\,),\quad \phii(x) = 6x+1,
\qquad
\psii : \Nsstar \xrightarrow{\;\sim\;} (6\Z-1,\,\bul),\quad \psii(x) = 6x-1
\]
are monoid isomorphisms.
\textnormal{Lean 4:} \href{\GH/TPC/Basic.lean\#L55}{\texttt{Basic.lean\#L55}}
\end{theorem}

\begin{proof}
$\phii(x \mstar y) = 6(6xy+x+y)+1 = (6x+1)(6y+1) = \phii(x)\phii(y)$.
$\psii(x \sstar y) = 6(-6xy+x+y)-1 = -[(6x-1)(6y-1)] = \psii(x) \bul \psii(y)$.
Both are bijections with inverses $n \mapsto (n-1)/6$ and $n \mapsto (n+1)/6$.
\end{proof}

\begin{lemma}[No zero divisors]
\label{lem:nzd}
If $u, v \neq 0$ then $u \mstar v \neq 0$ and $u \sstar v \neq 0$.
\textnormal{Lean 4:} \href{\GH/TPC/Basic.lean\#L90}{\texttt{Basic.lean\#L90}}
\end{lemma}

\begin{proof}
$u \mstar v = 0 \iff (6u+1)(6v+1) = 1$ in $\Z$, forcing $6u+1 = \pm 1$,
giving $u = 0$ or $u = -1/3 \notin \Z$.  Similarly for $\sstar$.
\end{proof}

\begin{definition}[Product monoid]
The \textbf{product monoid} is $M = (\Z^2, \tens)$ where
\[
(x,y) \tens (x',y') = (x \mstar x',\; y \sstar y'), \qquad \mathbf{1}_M = (0,0).
\]
\textnormal{Lean 4:} \href{\GH/TPC/Monoid.lean\#L40}{\texttt{Monoid.lean\#L40}}
\end{definition}

% ═════════════════════════════════════════════════════════════════════════
\section{The Anti-Diagonal Submonoid}
% ═════════════════════════════════════════════════════════════════════════

\begin{definition}[Anti-diagonal]
$N = \Deltaop = \{(x,-x) : x \in \Z\} \subset M$.
\textnormal{Lean 4:} \href{\GH/TPC/AntiDiagonal.lean\#L10}{\texttt{AntiDiagonal.lean\#L10}}
\end{definition}

\begin{proposition}
\label{prop:submonoid}
$N$ is a submonoid of $M$, isomorphic to $\Nstar$ via $\phii_\Delta(k)=(k,-k)$.
\[
\phii_\Delta(x) \tens \phii_\Delta(y) = \phii_\Delta(x \mstar y)
\qquad \text{(\texttt{Lean}: \texttt{Δ\_otimes})}.
\]
\end{proposition}

\begin{proof}
$(-x) \sstar (-y) = -6xy-x-y = -(x \mstar y)$, so
$(x,-x) \tens (y,-y) = (x \mstar y, -(x \mstar y)) \in N$.
\end{proof}

\noindent\textbf{The form set and twin primes.}\quad
Define $F = \{6xy+x+y : x,y \neq 0\} \cup \{-6xy+x+y : x,y \neq 0\}$.
Via $\phii$ and $\psii$:
\[
n \notin F \iff 6n+1 \in \pm\PP \text{ and } 6n-1 \in \pm\PP,
\]
i.e.\ $(6n-1, 6n+1)$ is a twin prime pair.  Therefore:
\[
\text{Twin Prime Conjecture} \iff |F^c| = \infty \iff |\Irr(\Nstar) \cap \Irr(\Nsstar)| = \infty.
\]

% ═════════════════════════════════════════════════════════════════════════
\section{Weak Saturation of $N$}
% ═════════════════════════════════════════════════════════════════════════

\noindent\textbf{Counterexample: $N$ is not saturated.}\quad
Take $a=(1,0)$, $b=(1,-8)$:
$(1,0) \tens (1,-8) = (1 \mstar 1, 0 \sstar (-8)) = (8,-8) \in N$,
but $(1,0) \notin N$ since $0 \neq -1$.
\textnormal{Lean 4:} \href{\GH/TPC/WeaklySaturated.lean\#L80}{\texttt{WeaklySaturated.lean\#L80}}

\begin{theorem}[$N$ is weakly saturated]
\label{thm:ws}
\textnormal{Lean 4:} \href{\GH/TPC/WeaklySaturated.lean\#L45}{\texttt{WeaklySaturated.lean\#L45}}
\end{theorem}

\begin{proof}
Suppose $(k,-k) = (a,b) \tens (c,d) \in N$ with $(a,b),(c,d) \neq (0,0)$.
Then $a \mstar c = k \neq 0$, so Lemma~\ref{lem:nzd} gives $a, c \neq 0$.
Set $a' = (a,-a) \in N$ and $b' = (c,-c) \in N$.  Then
\[
a' \tens b' = (a \mstar c,\; (-a) \sstar (-c)) = (k, -(a \mstar c)) = (k,-k),
\]
using $(-a) \sstar (-c) = -(a \mstar c)$ by the isomorphism $\etaa$.
This exhibits a nontrivial factorization within $N$.
\end{proof}

\begin{example}
The counterexample lifts: $(1,0) \tens (1,-8) = (8,-8)$ becomes
$(1,-1) \tens (1,-1) = (8,-8)$ within $N$.
\end{example}

% ═════════════════════════════════════════════════════════════════════════
\section{Irreducibility in $N$}
% ═════════════════════════════════════════════════════════════════════════

\begin{theorem}[Irreducibility equivalence]
\label{thm:irred-equiv}
For $(x,-x) \in N$ with $x \neq 0$:
\[
(x,-x) \in \Irr(N) \iff x \in \Irr(\Nstar) \cap \Irr(\Nsstar)
\iff 6x+1 \in \pm\PP \text{ and } 6x-1 \in \pm\PP.
\]
\textnormal{Lean 4:} \href{\GH/TPC/AntiDiagonal.lean\#L95}{\texttt{AntiDiagonal.lean\#L95}}
\end{theorem}

\begin{proof}
$(\Leftarrow)$: If $(x,-x) = (a,-a) \tens (b,-b)$ nontrivially, the first
component gives $a \mstar b = x$ nontrivially, contradicting
$x \in \Irr(\Nstar)$.

$(\Rightarrow)$: If $x = a \mstar b$ nontrivially, then
$(a,-a) \tens (b,-b) = (x,-x)$ is nontrivial in $N$, contradiction.
For $\sstar$-irreducibility: if $x = a \sstar b$ nontrivially then
$\etaa$ gives a nontrivial factorization of $(x,-x)$ in $N$.
\end{proof}

\begin{center}
\begin{tabular}{rrrll}
\toprule
$x$ & $6x-1$ & $6x+1$ & $\Irr(N)$? & Twin pair \\
\midrule
$1$  & $5$  & $7$  & Yes & $(5,7)$ \\
$2$  & $11$ & $13$ & Yes & $(11,13)$ \\
$3$  & $17$ & $19$ & Yes & $(17,19)$ \\
$4$  & $23$ & $25=5^2$ & No & --- \\
$5$  & $29$ & $31$ & Yes & $(29,31)$ \\
$7$  & $41$ & $43$ & Yes & $(41,43)$ \\
$10$ & $59$ & $61$ & Yes & $(59,61)$ \\
\bottomrule
\end{tabular}
\end{center}

% ═════════════════════════════════════════════════════════════════════════
\section{The Norm}
% ═════════════════════════════════════════════════════════════════════════

\begin{definition}[The norm]
\label{def:norm}
For $(x,-x) \in N$ define
\[
|(x,-x)| = \phii(x)^2 = (6x+1)^2.
\]
\textnormal{Lean 4:} \href{\GH/TPC/Norm.lean\#L12}{\texttt{Norm.lean\#L12}}
\end{definition}

\begin{proposition}[Norm axioms]
The norm satisfies:
\begin{enumerate}[label=\textbf{N\arabic*.},leftmargin=*]
\item \textbf{Non-negativity.} $|(x,-x)| \geq 0$.
\item \textbf{Identity.} $|(0,0)| = 1$.
\item \textbf{Non-degeneracy.} $|(x,-x)| = 1 \iff x = 0 \iff (x,-x) = 1_N$.
\item \textbf{Integrality.} $|(x,-x)| \in \Z_{>0}$.
\item \textbf{Multiplicativity.}
  $|(x,-x) \tens (y,-y)| = |(x,-x)| \cdot |(y,-y)|$.
\item \textbf{Strict descent.} If $(x,-x) = (a,-a) \tens (b,-b)$ with
  $a,b \neq 0$ then $|(a,-a)|, |(b,-b)| < |(x,-x)|$.
\end{enumerate}
\end{proposition}

\begin{proof}
N5: $|(x \mstar y, -(x \mstar y))| = \phii(x \mstar y)^2 = (\phii(x)\phii(y))^2 = \phii(x)^2 \cdot \phii(y)^2$.
N6: from N5 and $|(a,-a)|, |(b,-b)| \geq 4$ (since $|6x+1| \geq 2$ for $x \neq 0$).
\end{proof}

\begin{center}
\begin{tabular}{rrrll}
\toprule
$x$ & $6x+1$ & $|(x,-x)|$ & $\Irr(N)$? & Note \\
\midrule
$0$  & $1$   & $1$    & Id  & identity \\
$1$  & $7$   & $49$   & Yes & $(5,7)$ twin \\
$2$  & $13$  & $169$  & Yes & $(11,13)$ twin \\
$3$  & $19$  & $361$  & Yes & $(17,19)$ twin \\
$4$  & $25$  & $625$  & No  & $25=5^2$ \\
$-1$ & $-5$  & $25$   & Yes & $(-7,-5)$ twin \\
\bottomrule
\end{tabular}
\end{center}

% ═════════════════════════════════════════════════════════════════════════
\section{Every Element has an Irreducible Divisor}
% ═════════════════════════════════════════════════════════════════════════

\begin{definition}
$(q,-q) \mid_\tens (k,-k)$ if $\exists (m,-m) \in N$ with
$(q,-q) \tens (m,-m) = (k,-k)$.
\end{definition}

\begin{theorem}
\label{thm:irred-div}
For every $(k,-k) \in N \setminus \{1_N\}$ there exists
$(q,-q) \in \Irr(N)$ with $(q,-q) \mid_\tens (k,-k)$.
\textnormal{Lean 4:} \href{\GH/TPC/Norm.lean\#L90}{\texttt{Norm.lean\#L90}}
\end{theorem}

\begin{proof}
Strong induction on $|(k,-k)| \geq 4$.  If $(k,-k)$ is irreducible: done.
Otherwise it is reducible in $M$; by Theorem~\ref{thm:ws} lift to
$(k,-k) = (a,-a) \tens (b,-b)$ nontrivially in $N$.  By N6,
$|(a,-a)| < |(k,-k)|$; induction gives $(q,-q) \in \Irr(N)$ with
$(q,-q) \mid_\tens (a,-a)$, hence $(q,-q) \mid_\tens (k,-k)$.
\end{proof}

\begin{corollary}[Covering]
\label{cor:cover}
$N \setminus \{(0,0)\} = \bigcup_{q \in \Irr(N)} (q,-q) \tens N$.
\end{corollary}

% ═════════════════════════════════════════════════════════════════════════
\section{The Topology}
% ═════════════════════════════════════════════════════════════════════════

\begin{definition}[Basis elements]
\label{def:U}
For $a,b,c \in \Z$, $c \neq 0$:
\[
U_{a,b,c} = \{(a+cz,\; b-cz) : z \in \Z\}.
\]
Every $(x,y) \in U_{a,b,c}$ satisfies $x+y = a+b$.
\textnormal{Lean 4:} \href{\GH/TPC/Topology.lean\#L18}{\texttt{Topology.lean\#L18}}
\end{definition}

\begin{example}
$U_{1,-1,7} = \{\ldots,(-6,6),(1,-1),(8,-8),(15,-15),\ldots\}$ ---
the arithmetic progression of step $7$ through $(1,-1)$ on $N$.
$U_{0,0,5} \cap N = \{(5n,-5n) : n \in \Z\}$ --- multiples of $5$ in $N$.
\end{example}

\begin{theorem}[$\mathcal{B}$ is a topological basis; each $U_{a,b,c}$ is clopen]
\label{thm:basis}
$\mathcal{B} = \{U_{a,b,c} : a,b,c \in \Z,\, c \neq 0\}$ is a basis for a
topology $\tau$ on $\Z^2$, and every $U_{a,b,c}$ is clopen.
\textnormal{Lean 4:} \href{\GH/TPC/Topology.lean\#L75}{\texttt{Topology.lean\#L75}}
\end{theorem}

\begin{proof}
\textbf{B1:} $(x,y) \in U_{x,y,1}$.
\textbf{B2 (intersections):} Reduce to
$x \equiv a \pmod{c}$, $x \equiv a' \pmod{c'}$; CRT gives
$U_{a,b,c} \cap U_{a',b',c'} = U_{x^*,\,a+b-x^*,\,\mathrm{lcm}(c,c')} \in \mathcal{B}$
when $a+b = a'+b'$.
\textbf{Clopen:} complement is
$\bigcup_{k=1}^{|c|-1} U_{a+k,b-k,c}$, a finite union of opens.
\end{proof}

\begin{theorem}[Subspace topology on $N$]
\label{thm:subspace}
\[
N \cap U_{a,b,c} =
\begin{cases}
\emptyset & \text{if } c \nmid a+b, \\
\{(a+cn,-(a+cn)) : n \in \Z\} & \text{if } c \mid a+b.
\end{cases}
\]
In particular every nonempty open set in $N$ is infinite; no singleton is open.
\end{theorem}

\begin{proof}
$(z,-z) \in U_{a,b,c}$ requires $0 = a+b$ (add the two coordinate conditions),
so $c \mid a+b$ is necessary; when satisfied $z \in a+c\Z$ is infinite.
\end{proof}

\begin{theorem}[$N$ is a topological monoid; $L_p$ is a homeomorphism]
\label{thm:topmonoid}
$(\Deltaop, \tens)$ is a topological monoid.  For
$p = (p_1,-p_1) \in N \setminus \{(0,0)\}$, left multiplication
$L_p(x,-x) = p \tens (x,-x)$ is a homeomorphism onto
$p \tens N = V_{p_1,\,6p_1+1}$.  In particular, for any closed
$C \subseteq N$ the set $p \tens C$ is closed in $N$.
\textnormal{Lean 4:} \href{\GH/TPC/Topology.lean\#L155}{\texttt{Topology.lean\#L155}}
\end{theorem}

\begin{proof}
Continuity: for $x = x_0 + cs$ and $y = y_0 + ct$,
$x \mstar y \equiv x_0 \mstar y_0 \pmod{c}$ since correction terms are
divisible by $c$.  Injectivity: $(6p_1+1)(6x+1) = (6p_1+1)(6y+1)$ and
$6p_1+1 \neq 0$ give $x = y$.  Open onto image: $L_p(V_{x_0,c}) =
V_{p_1 \mstar x_0,\,(6p_1+1)c}$, a basis element.  Closure follows because
$L_p(\Deltaop) = V_{p_1,6p_1+1}$ is clopen, so $L_p(C)$ is closed in it.
\end{proof}

% ═════════════════════════════════════════════════════════════════════════
\section{Main Theorem}
% ═════════════════════════════════════════════════════════════════════════

\begin{mainresult}
\begin{theorem}[Main theorem]
\label{thm:main}
$|\Irr(N)| = \infty$.
\textnormal{Lean 4:} \href{\GH/TPC/Main.lean\#L60}{\texttt{Main.lean\#L60}}
\end{theorem}
\end{mainresult}

\begin{proof}
By Corollary~\ref{cor:cover}:
\[
N \setminus \{(0,0)\} = \bigcup_{q \in \Irr(N)} q \tens N.
\]
Suppose $\Irr(N) = \{q^{(1)}, \ldots, q^{(n)}\}$ is finite.  By
Theorem~\ref{thm:topmonoid} each $q^{(i)} \tens N$ is closed.  A finite
union of closed sets is closed, so $N \setminus \{(0,0)\}$ is closed.
Therefore $\{(0,0)\}$ is open in $N$.  But by Theorem~\ref{thm:subspace}
every nonempty open set is infinite --- a singleton cannot be open.
Contradiction.
\end{proof}

\begin{remark}[Furstenberg analogy]
This argument is verbatim Furstenberg's 1955 proof~\cite{Furstenberg1955}:
$\Z \setminus \{-1,1\} = \bigcup_p p\Z$ (closed); finiteness forces
$\{-1,1\}$ open; contradiction.  Here $N$ plays the role of $\Z$ and
$U_{a,b,c}$ plays the role of $a + b\Z$.
\end{remark}

\begin{mainresult}
\begin{corollary}[Twin Prime Conjecture]
\label{cor:tpc}
\[
\bigl|\{\,x \in \Z : 6x-1 \in \pm\PP \text{ and } 6x+1 \in \pm\PP\,\}\bigr| = \infty.
\]
That is, there are \textbf{infinitely many twin prime pairs}.
\end{corollary}
\end{mainresult}

\begin{proof}
Immediate from Theorems~\ref{thm:irred-equiv} and~\ref{thm:main}.
\end{proof}

% ═════════════════════════════════════════════════════════════════════════
\section*{Acknowledgements}
% ═════════════════════════════════════════════════════════════════════════

Computational and expository assistance from Claude Sonnet 4.6 (Anthropic).
The Lean 4 formalisation uses Mathlib4.

% ─────────────────────────────────────────────────────────────────────────
\begin{thebibliography}{9}
\bibitem{Furstenberg1955}
H.~Furstenberg,
\emph{On the infinitude of primes},
Amer.\ Math.\ Monthly \textbf{62} (1955), 353.

\bibitem{AtiyahMacdonald}
M.~F.~Atiyah and I.~G.~Macdonald,
\emph{Introduction to Commutative Algebra},
Addison-Wesley, 1969.

\bibitem{HardyWright}
G.~H.~Hardy and E.~M.~Wright,
\emph{An Introduction to the Theory of Numbers}, 6th ed.,
Oxford University Press, 2008.

\bibitem{Serre}
J.-P.~Serre,
\emph{A Course in Arithmetic},
Springer-Verlag, 1973.

\bibitem{Mathlib4}
The Mathlib4 Community,
\emph{Mathlib4},
\url{https://leanprover-community.github.io/mathlib4_docs/}, 2024.
\end{thebibliography}

\end{document}
